% ===================================================================
%  图表第 14 号 — 街道。— 商人。
%  Traduction 2026 d'apres texto/io, controlee sur texto/fr, texto/en
%  et texto/es. Consigne : voir texto/zh/10-tubiao-01.tex.
% ===================================================================

%%K t14-apar-1 apar
\VUsaut{4.0mm}
\VUtitre{15.0pt}{60}{图表第 14 号}
\VUsaut{3.0mm}

%%K t14-tit-1 sub
\VUpk{11.4pt}{40}{街道。--- 商人。}
\VUsaut{3.0mm}

%%K t14-01-1 p
1. --- 我的朋友约翰住在乡下，来我家住三天。\VUgras{直到}如今他还没有机会
见过大城市，所以我们整天在外面走，他不懂的我就讲给他听。

%%K t14-02-1 p
2. --- 我们先去看\VUgras{天文台} \textsuperscript{(1)}，他很感兴趣。回来时走那条\VUgras{街}
\textsuperscript{(3)}，就是有\VUgras{土耳其浴室} \textsuperscript{(2)} 的那条，碰上
一队\VUgras{出殡} \textsuperscript{(4)}。\VUgras{送葬队伍}最前头走着\VUgras{神职
人员}，在\VUgras{灵车}前面，\VUgras{棺材}就放在灵车上。许多朋友跟在
戴\VUgras{孝}的家属后面。

%%K t14-02-2 p
队伍过去以后，我们在大\VUgras{啤酒店} \textsuperscript{(5)} 前面穿过街去；许多
\VUgras{客人}在露台上喝\VUgras{啤酒}，头上有玻璃\VUgras{雨篷} \textsuperscript{(6)}
遮着。

%%K t14-03-1 p
3. --- \VUgras{纹章}挂在\VUgras{酒商} \textsuperscript{(7)} 的铺子和\VUgras{乐谱商}
\textsuperscript{(10)} 的铺子中间，约翰弄不明白。他问我那是什么意思；我告诉他
那是英国\VUgras{领事馆} \textsuperscript{(8)}，还指给他看站在阳台上的\VUgras{领事}
\textsuperscript{(9)}。

%%K t14-tit-2 sub
\VUpk{10.2pt}{40}{看橱窗。}
\VUsaut{3.0mm}

%%K t14-04-1 p
4. --- 我们自然在陈列\VUgras{乐器}、\VUgras{簧风琴}、\VUgras{管风琴}和
\VUgras{钢琴}的地方停了下来；我们看了\VUgras{乐谱}和钢琴\VUgras{曲谱}的
彩色封面，也赞叹那把当作招牌的镀金木\VUgras{竖琴}。

%%K t14-05-1 p
5. --- \VUgras{清道夫} \textsuperscript{(11)} 和\VUgras{扫街车} \textsuperscript{(12)}
扬起的尘土叫我们只好快步走过\VUgras{家具商} \textsuperscript{(13)} 那些漂亮的
\VUgras{成套家具}，也快步走过\VUgras{五金商} \textsuperscript{(14)} 陈列的
\VUgras{炉子}和\VUgras{灯}。

%%K t14-06-1 p
6. --- 就在那地方，我们真的不得不下人行道，因为道上堆着从\VUgras{搬家车}
\textsuperscript{(16)} 上卸下来、要搬进一间刚租出去的\VUgras{铺子} \textsuperscript{(15)}
的包裹。

%%K t14-06-2 p
这叫我们没看成\VUgras{马车匠} \textsuperscript{(17)} 那些漂亮车子，可是没耽误
我们停下来看\VUgras{包装工} \textsuperscript{(19)} 给他的邻居——那位
\VUgras{自行车和汽车商} \textsuperscript{(18)}——钉箱子。后来天色不早了，我们
就回家了。

%%K t14-tit-3 sub
\VUpk{10.2pt}{40}{在城里走一趟。}
\VUsaut{3.0mm}

%%K t14-07-1 p
7. --- 第二天母亲提议让约翰照张相，叫我陪他去\VUgras{照相馆} \textsuperscript{(20)}。
午饭后我们去了那位艺术家的\VUgras{照相室}，在那里等了好一会儿。为了消磨
时间，我们看墙上挂的\VUgras{相片} \textsuperscript{(22)}。

%%K t14-07-2 p
轮到我们时，事情没费多少工夫；我朋友一在\VUgras{照相机} \textsuperscript{(21)}
前面摆好姿势照完，我们就下楼了。

%%K t14-08-1 p
8. --- 在二层，我们从\VUgras{理发店} \textsuperscript{(23)} 敞开的门里看见伙计
正给一位客人剪头。

%%K t14-08-2 p
又回到街上，我们经过\VUgras{烟草店} \textsuperscript{(24)}。约翰给那盏红
\VUgras{灯} \textsuperscript{(25)} 和当\VUgras{招牌} \textsuperscript{(26)} 的
\VUgras{大烟斗}逗得直乐，一头撞上了两位正一面闻\VUgras{鼻烟} \textsuperscript{(27)}
一面闲谈的老先生。

%%K t14-tit-4 sub
\VUpk{10.2pt}{40}{点心铺。}
\VUsaut{3.0mm}

%%K t14-09-1 p
9. --- \VUgras{兑换商} \textsuperscript{(29)} \VUgras{橱窗}里陈列的各国
\VUgras{钞票}和\VUgras{硬币}把我们留住了一会儿，可是那时正是吃点心的时候，
我们没再多停就进了\VUgras{点心铺} \textsuperscript{(31)}。

%%K t14-10-1 p
10. --- 看见铺子里那么多\VUgras{好东西}——\VUgras{糕点、姜饼、饼干}和各样
\VUgras{糖果}——我们简直挑不出来。末了约翰选了\VUgras{奶油糕}，我只要了
一小块\VUgras{樱桃挞}，因为甜的东西\VUgras{叫我牙疼}，我一点也不想常去
\VUgras{牙医} \textsuperscript{(30)} 那里。

%%K t14-tit-5 sub
\VUpk{10.2pt}{40}{打字。}
\VUsaut{3.0mm}

%%K t14-11-1 p
11. --- 为了给朋友看一台\VUgras{打字机}——他听说过却没见过——我们走进
那间\VUgras{办公室} \textsuperscript{(32)}，是我\VUgras{表哥}
\textsuperscript{(34)} 的。我们
撞见他正在给他的\VUgras{女秘书} \textsuperscript{(35)} 口授信件，那位女秘书是
\VUgras{速记员}。一封信刚说完，一个\VUgras{打字员} \textsuperscript{(33)} 就在
机器上把它打出来。我们不愿打扰亲戚做事，只待了几分钟。

%%K t14-12-1 p
12. --- 我表哥的办公室底下住着一位大\VUgras{钟表匠兼首饰商} \textsuperscript{(37)}，
他还是个银匠。他那漂亮铺子里有许多\VUgras{挂钟、座钟、怀表}、
\VUgras{表链}和贵重的\VUgras{首饰}，像\VUgras{珍珠项链}、金\VUgras{手镯}、
\VUgras{耳环}和嵌着\VUgras{宝石}和\VUgras{钻石}的\VUgras{戒指}。有一个橱窗
整个给了\VUgras{银器}，里面有一大批银\VUgras{餐具}。

%%K t14-13-1 p
13. --- 拐过街角时，我注意到房子\VUgras{门牌}旁边那块\VUgras{牌子}
\textsuperscript{(36)} 换过了。我正要说出口，就听见军乐队欢快的乐声。我们朝那
队伍跑去，没想到它要经过\VUgras{帽商} \textsuperscript{(39)} 和\VUgras{手套商}
\textsuperscript{(40)} 的铺子；其实我们要是待在\VUgras{眼镜商} \textsuperscript{(38)}
的橱窗前面看队伍走过，位置一样好，那橱窗里满是\VUgras{夹鼻眼镜、眼镜}和
\VUgras{观剧镜}。

%%K t14-tit-6 sub
\VUpk{10.2pt}{40}{队伍走过。}
\VUsaut{3.0mm}

%%K t14-14-1 p
14. --- \VUgras{团} \textsuperscript{(41)} 的最前头走着\VUgras{鼓手长}
\textsuperscript{(42)}，后面是\VUgras{鼓手} \textsuperscript{(43)}、\VUgras{号手}
\textsuperscript{(44)} 和整个\VUgras{乐队}。\VUgras{将军} \textsuperscript{(45)} 带着
副官在广场上，停在\VUgras{水柱} \textsuperscript{(49)} 旁边看
\VUgras{阅兵}，那水柱是\VUgras{喷泉} \textsuperscript{(48)} 的。

%%K t14-14-2 p
\VUgras{参谋军官} \textsuperscript{(46)}、\VUgras{上校}和\VUgras{中校}举剑向他
致敬。\VUgras{少校}们走在自己的\VUgras{营} \textsuperscript{(47)} 前头，每个
\VUgras{上尉}带着自己的\VUgras{连}，\VUgras{中尉}、\VUgras{少尉}和
\VUgras{士官}都在自己人旁边照常的位置上。

%%K t14-15-1 p
15. --- \VUgras{十字路口} \textsuperscript{(50)} 聚了许多行人；铺子里的人——
\VUgras{伞商} \textsuperscript{(51)}、\VUgras{染坊主} \textsuperscript{(53)}、
\VUgras{枪械商} \textsuperscript{(54)}——都出来看\VUgras{兵}。所有的车——
\VUgras{出租马车} \textsuperscript{(52)}、\VUgras{自用马车} \textsuperscript{(57)}，
连\VUgras{洒水车} \textsuperscript{(56)} 也在内——都靠边停下。

%%K t14-tit-7 sub
\VUpk{10.2pt}{40}{街头的景象。}
\VUsaut{3.0mm}

%%K t14-15-2 p
一个手提\VUgras{样品箱}的\VUgras{行商} \textsuperscript{(55)} 快步穿过街去，坐在\VUgras{顶层}
\textsuperscript{(59)} 上的人都回过头来看热闹，那是\VUgras{公共马车}
\textsuperscript{(58)} 的顶层。一个可怜的\VUgras{街头歌手} \textsuperscript{(60)} 带着他的
\VUgras{手摇风琴} \textsuperscript{(61)}，只好停下他的\VUgras{歌}，因为鼓声把他
的嗓子盖住了。

%%K t14-16-1 p
16. --- 队伍走到\VUgras{呢绒店} \textsuperscript{(64)} 前面时，\VUgras{女缝工}
\textsuperscript{(63)} 和\VUgras{裁缝} \textsuperscript{(62)} 丢下\VUgras{缝纫机}和
活计，跑到窗口去看。\VUgras{店主} \textsuperscript{(65)} 刚把新\VUgras{衣服}
\textsuperscript{(66)} 摆进\VUgras{橱窗}，只好把\VUgras{呢子} \textsuperscript{(67)} 和\VUgras{布}卷搬开——
那些卷摊在人行道上，靠近\VUgras{下水口} \textsuperscript{(68)}——
怕人群把它们碰倒；连一个老\VUgras{货郎} \textsuperscript{(69)} 也只好躲进门洞
去做完他的买卖，那时一个看门的女人正跟他讲袜子的价钱。

%%K t14-16-2 p
我们送那队伍一直到兵营，\VUgras{然后}很满意地在晚饭时候回了家。

%%K t14-tit-8 sub
\VUpk{10.2pt}{40}{各样的行业和职业。}
\VUsaut{3.0mm}

%%K t14-17-1 p
17. --- 第二天父亲提议带我们去看本地报纸的印刷所。我们高高兴兴地答应了。

%%K t14-17-2 p
那位\VUgras{印刷商}同时也是\VUgras{书商} \textsuperscript{(70)}、\VUgras{出版商}、
\VUgras{文具商}和\VUgras{装订商}。他把报纸的\VUgras{编辑部} \textsuperscript{(71)}
安在二层，还有\VUgras{制版间}，\VUgras{石印工} \textsuperscript{(72)} 和
\VUgras{铜版雕工} \textsuperscript{(73)} 在那里做活；最后是排字间，\VUgras{排字工}
\textsuperscript{(74)} 在那里。真正的\VUgras{机器房} \textsuperscript{(75)}、
\VUgras{印刷机}和各样机器都在底层。

%%K t14-tit-9 sub
\VUpk{10.2pt}{40}{约翰问个不停。}
\VUsaut{3.0mm}

%%K t14-18-1 p
18. --- 从印刷所出来，约翰很纳闷地停在\VUgras{绒线店} \textsuperscript{(77)}
对面一根柱子上的小玻璃匣前，问它做什么用。父亲告诉他那是\VUgras{火警器}
\textsuperscript{(76)}。的确，城里样样东西对我朋友都新鲜：从那简单的
\VUgras{饮水泉} \textsuperscript{(78)}、轻便的\VUgras{送货三轮} \textsuperscript{(80)}——
骑车的是穿号衣的\VUgras{侍应生} \textsuperscript{(79)}——一直到\VUgras{邮车}
\textsuperscript{(81)}。

%%K t14-19-1 p
19. --- 父亲离开我们一会儿去同\VUgras{税吏} \textsuperscript{(83)} 说话——那人
正停下同一位\VUgras{律师} \textsuperscript{(82)} 和一位\VUgras{公证人}
\textsuperscript{(84)} 谈天——我们就看街上的来往解闷。各色人等从我们面前过去：
一个\VUgras{点心铺学徒} \textsuperscript{(85)} 用篮子提着一个漂亮的\VUgras{馅饼}，
跟在一个\VUgras{估衣商} \textsuperscript{(86)} 后面；一个\VUgras{点灯人}
\textsuperscript{(87)} 在同一个\VUgras{煤气工} \textsuperscript{(88)} 说话；一位
\VUgras{修女} \textsuperscript{(89)} 牵着一个小孤女的手往修道院走。

%%K t14-tit-10 sub
\VUpk{10.2pt}{40}{回家路上。}
\VUsaut{3.0mm}

%%K t14-20-1 p
20. --- 父亲刚回到我们身边，约翰就大笑起来，因为看见两个\VUgras{扫烟囱的}
\textsuperscript{(91)} 满脸黑灰、打扮古怪。

%%K t14-21-1 p
21. --- 我想向\VUgras{卖花的女人} \textsuperscript{(92)} 给母亲买一盆花，可是父亲说
那太沉，不好拿，不如买些\VUgras{橘子}。我们走到\VUgras{摊贩} \textsuperscript{(94)}
跟前；等那位正给孩子挑橘子的\VUgras{家庭女教师} \textsuperscript{(93)} 买完，
我们才买上。

%%K t14-22-1 p
22. --- 回家路上我们遇见好几个熟人：我家一位老朋友挽着她的\VUgras{女伴}
\textsuperscript{(95)} 的胳膊；还有路桥\VUgras{工程师} \textsuperscript{(96)}；还有我
一位表姐带着她的\VUgras{保姆} \textsuperscript{(98)}。

%%K t14-22-2 p
后来我们又走过我母亲的\VUgras{女帽商} \textsuperscript{(97)} 那里，还遇见两个
外乡的\VUgras{修士} \textsuperscript{(99)}，他们过来问我父亲去车站的路。
\VUsaut{6.0mm}
\VUfilet{22mm}
